\chapter{Metodologia de Pesquisa}
\label{pre:metodologia}

A abordagem de pesquisa proposta neste pr\'e-projeto envolve um conjunto de etapas planejadas para investigar como redes neurais de grafos podem ser aplicadas ao reconhecimento de padr\~oes em dados educacionais.

\section{Coleta de dados}

\begin{enumerate}
    \item Sele\c{c}\~ao de fontes de dados educacionais, como registros acad\^emicos, informa\c{c}\~oes de desempenho, dados demogr\'aficos e outras vari\'aveis relevantes.
    \item Obten\c{c}\~ao de permiss\~ao de acesso aos dados, com aten\c{c}\~ao \`as quest\~oes de privacidade e \'etica.
\end{enumerate}

\section{Pr\'e-processamento}

\begin{enumerate}
    \item Limpeza de dados para remo\c{c}\~ao de valores faltantes, inconsist\^encias e \textit{outliers}.
    \item Transforma\c{c}\~ao dos dados conforme necess\'ario, incluindo padroniza\c{c}\~ao, codifica\c{c}\~ao de vari\'aveis categ\'oricas e demais ajustes pertinentes.
\end{enumerate}

\section{Modelagem com redes neurais de grafos}

\begin{enumerate}
    \item Sele\c{c}\~ao da arquitetura apropriada para a natureza dos dados, como \sigla{GCN}{Graph Convolutional Network} ou \sigla{GNN}{Graph Neural Network}.
    \item Treinamento do modelo com ajuste de hiperpar\^ametros e monitoramento por m\'etricas adequadas.
    \item Valida\c{c}\~ao cruzada para avaliar a capacidade de generaliza\c{c}\~ao para novos dados.
\end{enumerate}

\section{Desenvolvimento do modelo}

\begin{enumerate}
    \item Defini\c{c}\~ao do objetivo do modelo como previs\~ao de desempenho e identifica\c{c}\~ao de problemas acad\^emicos.
    \item Sele\c{c}\~ao de vari\'aveis de entrada relevantes, como notas, desempenho acad\^emico, dados demogr\'aficos e outras vari\'aveis.
    \item Escolha da arquitetura de rede neural convolucional de grafo.
    \item Treinamento e avalia\c{c}\~ao do modelo nos dados educacionais.
\end{enumerate}

\section{An\'alise dos resultados}

\begin{enumerate}
    \item Avalia\c{c}\~ao do desempenho em termos de precis\~ao, sensibilidade, especificidade e demais m\'etricas adequadas.
    \item Interpreta\c{c}\~ao dos padr\~oes identificados pelo modelo para extrair \textit{insights} educacionais significativos.
\end{enumerate}

\section{Limita\c{c}\~oes e desafios}

O estudo prev\^e a identifica\c{c}\~ao e discuss\~ao de limita\c{c}\~oes, como poss\'iveis distor\c{c}\~oes nos dados, quest\~oes de interpretabilidade do modelo e restri\c{c}\~oes \'eticas.

\section{Conclus\~oes e recomenda\c{c}\~oes futuras}

Ao final, o trabalho dever\'a sintetizar os principais resultados e contribui\c{c}\~oes e indicar caminhos para pesquisas futuras que ampliem ou aprimorem a aplica\c{c}\~ao de redes neurais de grafos a dados educacionais.
